\documentclass[11pt,a4paper]{article}

\usepackage[LGR,T1]{fontenc}
\usepackage[utf8]{inputenc}
\usepackage[greek.ancient,spanish]{babel}
\usepackage[a4paper,margin=2.8cm]{geometry}
\usepackage{microtype}
\usepackage[hidelinks]{hyperref}

% Griego politónico escrito directamente en UTF-8:
% \gr{...} para palabras sueltas; bloques con \foreignlanguage{greek}{...}
\newcommand{\gr}[1]{\foreignlanguage{greek}{#1}}

\hypersetup{
  pdftitle={Leer la Odisea en tiempos iletrados. Segunda parte: Una osada Odisea}
}

\setlength{\parindent}{1.25em}
\setlength{\parskip}{0.35em}
\raggedbottom

\title{Leer la \emph{Odisea} en tiempos iletrados \\ Segunda parte \\ Una osada Odisea}
\author{}
\date{}

\begin{document}

\maketitle

Hace un par de semanas, mi hijo Héctor y yo nos fuimos de compras.

Primero fuimos a Zara a comprar unas camisas, es decir, a que él comprara camisas para ambos. Luego hicimos algunos otros recados, preparando su inminente viaje en tren por Europa, el interrail que une los veranos de tantos jóvenes del continente. Una parte importante de su itinerario pasaba por Grecia, desde donde me mandó ayer unas fotos, sentado con un amigo junto al Ponto de color de vino, justo enfrente de la isla de Ítaca. Confieso sin vergüenza que las fotos me hicieron llorar. Allí estaba mi hijo, tan noble y hermoso como aquel otro hombre de la leyenda, allí estaba un sol color escarlata iluminando las costas de un lugar cuyo nombre equivale a todo lo que nos es querido, a todo lo que nos llama, a todo lo que hemos perdido.

Ítaca.

Padre, hijo, hogar, esposa, familia. La historia de la \emph{Odisea} no puede ser más familiar ni se puede haber contado más veces. Achab, Nemo, Simbad, incluso el mismísimo John Silver, son trasuntos de Ulises. Borges aseguraba que toda la historia de la literatura occidental se reducía a cuatro relatos ---una ciudad asaltada a traición, un marino que intenta regresar a su hogar, una búsqueda, un hombre que muere en la cruz--- y quizás no anduviera desencaminado.

La última tienda que visitamos Héctor y yo fue una librería, ya que tenía el secreto propósito de hacerle el mismo regalo ---la \emph{Ilíada} y la \emph{Odisea}--- que me hizo mi padre a mí hace medio siglo, para que le acompañaran en su viaje a Grecia. Y en efecto, gracias a Nolan, las estanterías rebosaban de ejemplares de los clásicos. Entre las versiones que estuve ojeando, todas ellas primorosamente encuadernadas y con bonitas ilustraciones en las tapas, estaba, cómo no, la venerable traducción en prosa de Segalá y varias traducciones en verso, incluyendo las ya clásicas de Fernando Gutiérrez (1953, Penguin) y José Manuel Pabón (1982, Gredos). Pero la que seleccioné para Héctor fue la versión de Juan Manuel Rodríguez Tobal, publicada muy recientemente (este mismo verano de 2026), en una cuidadísima edición bilingüe de Hiperión.

Puedo imaginar al agudo lector, a la inteligente lectora preguntándose por mis razones y la pregunta es todo menos baladí. ¿Cuáles son los criterios para escoger una traducción concreta cuando se trata de un clásico que nos ha acompañado durante tres mil años?

El propio Tobal escribe, en un brevísimo prólogo, refiriéndose a su nueva traducción:

\begin{quotation}
Y aquí he llegado. Y llegar aquí parece que es labor que exige algún argumento.
\end{quotation}

Ciertamente, Tobal, profesor de lenguas clásicas, editor, poeta y reconocidísimo traductor de Safo, Catulo, Ovidio y Anacreonte, entre otros, no se ha caído de ningún guindo. Conoce ---al dedillo--- todas las traducciones clásicas al castellano (y sin duda a otras lenguas) de la \emph{Odisea} y no ignora que atreverse con otra «exige algún argumento». Continúa así:

\begin{quotation}
Claro, la traducción que doy en estas páginas es una, esta sola, una composición particular de entre las infinitas posibilidades que tiene el poema de Homero de ser reproducido, producido de nuevo. Y cientos de personas, de diferentes épocas, edades, lenguas y civilizaciones, se entregaron antes que yo a esta misma labor.
\end{quotation}

Hay aquí un eco del célebre relato de Borges, «Pierre Menard, autor del Quijote», en el que un oscuro escritor francés decide ---y consigue--- reescribir varios capítulos del \emph{Quijote}, palabra por palabra, sin que estos sean una copia. Y para ello, naturalmente, no le queda otra que transformarse en Cervantes:

\begin{quotation}
[...] El método inicial que imaginó era relativamente sencillo. Conocer bien el español, recuperar la fe católica, guerrear contra los moros o contra el turco, olvidar la historia de Europa entre 1602 y 1918, ser Miguel de Cervantes. [...]
\end{quotation}

Escila y Caribdis, por supuesto. O bien aproximarse tanto al original como para confundirse con él ---y de paso hacerse incomprensible a los lectores del futuro, a la vez que irrelevante a los del presente, que ya disponen del original--- o bien interpretarlo libremente, corriendo el riesgo de contar, una vez más, una historia contada hasta el agotamiento. Hay que ser muy valiente o muy insensato para atreverse.

El valor, a los traductores, igual que a los militares, se les supone. Pero además, Tobal no parece especialmente insensato cuando escribe:

\begin{quotation}
¿Debería hacer yo un catálogo de las razones de por qué hice de esta manera lo que hice y no de esa otra? Si me lo impusiera, tendría que argumentar hasta el aburrimiento y la tristeza por qué a lo que tantos dijeron «no se puede» yo digo «sí se puede»: por qué el verso, por qué el hexámetro, por qué la cadencia musical (y no la cadencia de metrónomo), por qué la poesía, por qué la rima, por qué los neologismos, por qué los arcaísmos, por qué la diversidad dialectal, por qué otras lenguas, por qué los nombres parlantes, por qué nunca tan mismos los versos formulares, por qué y dónde y cómo la Ilíada... y muchos más porqués.
\end{quotation}

Tobal se niega al catálogo y al debate, pero ofrece al perspicaz lector, a la atrevida lectora, claves cruciales para orientarse en el océano que se dispone a abordar. Quizás podamos explorar algunas de ellas juntos.

Empezando por el principio, por el arranque de la \emph{Odisea} y la traducción clásica en prosa de Segalá:

\begin{quotation}
Háblame, Musa, de aquel varón de multiforme ingenio que, después de destruir la sacra ciudad de Troya, anduvo peregrinando larguísimo tiempo, vio las poblaciones y conoció las costumbres de muchos hombres y padeció en su ánimo gran número de trabajos en su navegación por el ponto, en cuanto procuraba salvar su vida y la vuelta de sus compañeros a la patria. Mas ni aun así pudo librarlos, como deseaba, y todos perecieron por sus propias locuras. ¡Insensatos! Comiéronse las vacas del Sol, hijo de Hiperión; el cual no permitió que les llegara el día del regreso. ¡Oh diosa, hija de Júpiter!: cuéntanos aunque no sea más que una parte de tales cosas.
\end{quotation}

Ya he contado, en la primera parte de este ensayo, que Segalá fue mi primera lectura de la \emph{Odisea}. Mirándolo en retrospectiva, me asombran las tragaderas que tenía mi yo adolescente. Hoy no sé si podría leer a don Luis con tanta felicidad como entonces y estoy convencido de que más de un atrevido lector, de una valerosa lectora, se darían de frente contra la pared de su espesa prosa. Por mi parte, no volví a leer el texto íntegro de la \emph{Odisea} en castellano, lo que no quiere decir que mi obsesión con Ulises se desvaneciera. Al contrario, fue a más, pero lo que me interesó, durante años, fue justo lo contrario de lo que ofrecía Segalá. Quería una reinterpretación moderna de la \emph{Odisea} y la encontré casi de carambola, cuando di con la magnífica \emph{Hija de Homero} de Robert Graves, una versión del mito que el escritor inglés compuso a partir de la tesis de Samuel Butler (\emph{The Authoress of the Odyssey}, 1897), según la cual la autora real del poema es una joven princesa de las colonias griegas de Sicilia llamada Nausícaa. La novela de Graves, de 1955, se adelanta en setenta años a la fiebre actual y maneja todo el feminismo inteligente que se le escapa a Nolan. Siguió la \emph{Odisea} de Kazantzakis (1938), en la que el héroe, aburrido en Ítaca, vuelve a partir y se enfrasca en mil nuevas aventuras y, cada cierto tiempo, alguna variante moderna, como \emph{The Lost Books of the Odyssey} (Zachary Mason, 2007). Es decir, después de asimilar la narración inicial y zamparme a Segalá, lo que me interesaba eran las vueltas de tuerca, la vista del mito desde el balcón contemporáneo. Para colmo, el \emph{Viatge a Ítaca} de Lluís Llach salió casi por la misma época, circa 1975. Llach me descubrió a Kavafis y Kavafis se convirtió en mi Homero particular y su Ítaca en la mía.

No volví a leer la \emph{Odisea} de un tirón hasta que Emily Wilson publicó su elegante traducción. Vale la pena ver cómo arranca:

\begin{verse}
Tell me about a complicated man.\\
Muse, tell me how he wandered and was lost\\
when he had wrecked the holy town of Troy,\\
and where he went, and who he met, the pain\\
he suffered on the sea, and how he worked\\
to save his life and bring his men back home.\\
He failed, and for their own mistakes, they died.\\
They ate the Sun God’s cattle, and the god\\
kept them from home. Now goddess, child of Zeus,\\
tell the old story for our modern times.\\
Find the beginning.
\end{verse}

Lo primero que noté al abrir el libro es que Wilson era concisa donde Segalá era engorroso y «moderna» donde don Luis sonaba arcaico. Asumí, por cierto, ignorante de mí, que Wilson usaba «verso libre». Nada más lejos de la realidad, pero permítame el culto lector, la docta lectora, alargar un poco más el suspense y citar ahora a Tobal:

\begin{verse}
De ese hombre, Musa, háblame, milerrante, que en tanto extravío,\\
dio tras dejar asolado de Troya el alcázar bendito,\\
Plazas de gentes y gentes él vio, y vio pensares distintos,\\
Y por el gran mar tragó en su alma dolores tantísimos,\\
Su aliento y vuelta amparando a los que llevaba consigo.\\
Pero por más que lo ansiara a salvar no llegó a sus amigos,\\
Que por sus propias locuras, las de ellos, cayeron perdidos,\\
¡Los insensatos! Las vacas del Sol, de Hiperión, Alto el Hijo,\\
Ellos comieron. Borró día en que uno regresa el dios mismo,\\
Hija de Zeus, de esto ya desde allá donde quieras tú dinos.
\end{verse}

Antes de seguir, inquieto lector, curiosa lectora, hazme un favor. Lee en voz alta, a ser posible declamando, aún mejor si lo haces en compañía y a orillas del mar. ¿Resuena en tus oídos la cadencia? ¿Sientes el ritmo?

Claro que sí y sin duda, siendo más astuto, más astuta que el autor de estas líneas, no pensarás que se trate de «verso libre». No lo es, ciertamente. Tobal escoge como base de su poema hexámetros castellanos. Cada hexámetro tiene seis apoyos o ictus. Se trata de la sílaba fuerte que da el latido, la cadencia del verso.

\begin{verse}
\emph{Hi}ja de \emph{Zeus}, de\_esto \emph{ya} desde\_a\emph{llá} donde \emph{quie}ras tú \emph{di}nos.
\end{verse}

Entre los seis apoyos (sílabas fuertes) del hexámetro (hi, Zeus, ya, llá, quie y di) hay cinco grupos (ja de, de\_esto, desde\_a, donde, ras tú) y un sexto tras el último apoyo (nos). Son los pies del hexámetro y los componen sílabas débiles. Dáctilo si contamos dos débiles, troqueo si contamos una. En el verso que nos ocupa, los cinco pies centrales son dáctilos (dos sílabas, incluyendo la sinalefa, que une vocales adyacentes) y el último pie es un troqueo. La regla general nos dice que el hexámetro puede usar dáctilo o troqueo y su composición «ideal» es que los dos últimos pies sean de la forma dáctilo-troqueo, tal como la línea que estamos analizando. El hexámetro de manual (como el que estamos viendo) tiene en total diecisiete sílabas (6 apoyos, 10 débiles en los cinco dáctilos y 1 en el troqueo final).

¿Son todos los hexámetros de Tobal de manual? Nada de eso. De hecho, abre con uno que no lo es:

\begin{verse}
(De\_ese) \emph{hom}bre \emph{Mu}sa \emph{há}blame\_ mile\emph{rran}te, que\_en \emph{tan}to\_ex tra\emph{ví}o,
\end{verse}

Aquí Tobal se salta el manual a la torera. En primer lugar, arranca con una anacrusa («De\_ese» son dos sílabas débiles) antes del primer apoyo en lugar de hacerlo en tiempo fuerte. A continuación comete impunemente un atropello de apoyos en el centro (los acentos de HOMbre, MUsa y HÁblame caen demasiado juntos: no deja dáctilo con cabeza, sustituyéndolos por troqueos). Arregla el desaguisado \emph{in extremis}, ya que en «e RRAN te que\_en TAN to\_ex tra VÍ o» la cosa se endereza y desemboca en la cláusula (los dos últimos pies, que son dáctilo y troqueo).

¿Hay un método en su locura? Yo diría que sí. La apertura sugiere que Tobal arranca el verso como le conviene (cediendo a la prosodia natural del español, que no siempre quiere empezar en tónica) y cierra siempre que puede en hexámetro de manual, inconfundiblemente dactílico. El conteo, si no me equivoco, sale así: empate entre anacrusas y tónicas en la cabeza del verso (cinco a cinco) y goleada del cierre canónico en la cola (siete a tres). Cuatro hexámetros regulares de diez y uno «perfecto».

Antes de seguir, echémosle un vistazo a la versión clásica de Pabón:

\begin{verse}
Musa, dime del hábil varón que en su largo extravío,\\
tras haber arrasado el alcázar sagrado de Troya,\\
conoció las ciudades y el genio de innúmeras gentes.\\
Muchos males pasó por las rutas marinas luchando\\
por sí mismo y su vida y la vuelta al hogar de sus hombres,\\
pero a éstos no pudo salvarlos con todo su empeño,\\
que en las propias locuras hallaron la muerte. ¡Insensatos!\\
Devoraron las vacas del Sol Hiperión e, irritada\\
la deidad, los privó de la luz del regreso. Principio\\
da a contar donde quieras, ¡oh diosa nacida de Zeus!
\end{verse}

Esta vez vamos a hacerlo al revés. Primero analizamos la estructura del fragmento, luego lo leemos en voz alta.

\begin{center}
\small
\begin{tabular}{cc|ccc|ccc|ccc|ccc|cc}
MU & sa & DI & me & del & HÁ & bil & va & RÓN & que\_en & su & LAR & go\_ex & tra & VÍ & o \\
X & -- & X & -- & -- & X & -- & -- & X & -- & -- & X & -- & -- & X & -- \\
\multicolumn{2}{c|}{T} & \multicolumn{3}{c|}{D} & \multicolumn{3}{c|}{D} & \multicolumn{3}{c|}{D} & \multicolumn{3}{c|}{D} & \multicolumn{2}{c}{T} \\
\end{tabular}
\end{center}

\begin{verse}
\emph{Mu}sa, \emph{di}me del \emph{há}bil va\emph{rón} que\_en su \emph{lar}go\_ex tra\emph{ví}o,
\end{verse}

Como en Tobal, seis apoyos, pero en este caso arrancando en tónica (MU-sa). Troqueo inicial, cinco dáctilos y troqueo final, con lo que la cláusula es canónica (dáctilo-troqueo). Si echamos cuentas, Pabón gestiona el arranque más o menos como Tobal (el conteo anacrusa a tónica sale 6 a 4, a comparar con 5 a 5) y se empeña en cerrar siempre canónico (diez de diez).

Y por supuesto, los dos intentan aproximarse al hexámetro dactílico del original, lo que no es nada fácil. El hexámetro griego es cuantitativo, no acentual. No mide por apoyos (sílabas tónicas), sino por duración, es decir, distingue entre sílabas largas (---) y breves (U). En otras palabras, en castellano el ictus que marca el ritmo viene dado por el acento. En el griego de la \emph{Odisea}, el ritmo viene marcado por el tiempo que se tarda en pronunciar la sílaba, algo impensable en castellano, que no distingue vocales largas de breves, pero también, irónicamente, en griego moderno, que tampoco las distingue. De ahí que la traducción canónica de la \emph{Odisea} al griego moderno, de Kazantzakis, recurra al verso de diecisiete sílabas, muy cercano, como vemos, a los hexámetros de Pabón y Tobal.

Entonces, ¿en qué se diferencian Pabón y Tobal? Hace falta leer más texto para entender la diferencia, pero, en última instancia, tiene que ver con el oído. Y aquí tengo que hacer un inciso para hablar de mi libro.



%Wilson me gusta mucho, es muy moderna, muy sutil y a ratos muy traditore. Toval me maravilla por su audacia. Con mi oído de cemento, jamás me habría atrevido (supongo que Jenaro sí) con una versión en hexámetros (lo más parecido que puede hacerse en castellano, me parece, al hexámetro dactílico del original).
%
%
%Para muestra de lo complicated que es la cosa (pun intended) Homero arranca así:
%
%ἄνδρα μοι ἔννεπε, μοῦσα, πολύτροπον,
%
%
%Literalmente: ἄνδρα (al hombre) — μοι (me) — ἔννεπε (refiere) — μοῦσα (Musa) πολύτροπον.
%
%πολύτροπον tiene miga. Poly —> mucho(s): tropon —> raíz tropos, con dos sentidos: hacer girar, volver, dar vueltas (de ahí milerrante en Toval) y por otra parte modo, manera, recurso… —> el de muchos recursos, o sea sagaz, astuto, de ahí Segalá (de multiforme ingenio). Y bastante espectacular Wilson con su “complicated”, que como ella misma explica tira de la raíz latina complicare, casi literalmente, “de muchos pliegues”, o “muchas vueltas”.
%
%Toval es el más fiel de los tres a la hora de arrancar: “De ese hombre, Musa, háblame, milerrante” es casi isomorfo al griego (excepto porque mil errante pierde el sentido de astuto). Segalá pone a la Musa delante del hombre (Háblame Musa de aquel varón) y se decanta por algo que se parece a la astucia (multiforme ingenio). Wilson también prioriza a la Musa (traditore!) y su “complicated” es a la vez una palabra “moderna” y un juego de palabras.
\end{document}
